% Results Section — v4 (2026-05-14)
% Changes from v3:
%   - Figure numbering: Fig 1 (paradigm), Fig 2 (forward encoder),
%                       Fig 3 (LORO/LOCO), Fig 4 (geometry),
%                       Fig 5 (2-component), Fig 6 (filter)
%   - Fig 1 (paradigm) moved to Methods to render before Fig 2 (forward encoder).
%   - Primary model: 2-component throughout; Machado and R+C comparisons moved to Appendix A
%   - LOCO metric updated: adjacent accuracy (adj_acc, 0–1, higher = better)
%     replacing circular correlation rho
%   - Fig 4 (SRM geometry/ΔRDM) added as separate figure
%   - Fig 1 is a PowerPoint-generated schematic: placeholder box below
%   - All captions from FIGURE_CAPTIONS.md (2026-05-11 revision)
% 2026-05-14 update (Option C adoption):
%   - Loss changed HYBRID → Option C (neural-primary 4-term + heavy Tikhonov, Methods Eq. 2).
%     Sub-08 argmin (16,+40) → (40,+26); sub-09 (12,-30) → (12,-28).
%   - Brettel directional-consistency text removed per 2026-05-14 revalidation.
%   - Sub-08 filter-direction description rewritten for new argmin.
%   - Model-class ceiling caveat added.
%   - HYBRID backup at results_v4_HYBRID_backup.tex.

\section{Results}
\label{sec:results}

% ──────────────────────────────────────────────────────────────────────────────
\subsection{Color discrimination is preserved across all ROIs in both CVD cases}
\label{sec:results:loro}
% ──────────────────────────────────────────────────────────────────────────────

A corrective stimulus-space filter is valid only if the visual cortex of a CVD individual retains categorical discrimination among the target hues.
We verified this precondition using leave-one-run-out (LORO) cross-validation of the forward encoding model in SRM-aligned space.

Both CVD participants exceeded the 0.125 exact-accuracy chance level at every ROI (Figure~\ref{fig:loco}A).
A cross-subject generalization test of the same forward encoding model (HC-to-HC vs.\ HC-to-CVD pair-wise cross-decoding, all ROIs pooled) showed no HC--CVD difference (Mann--Whitney $U$ test, 21 HC-to-HC vs.\ 14 HC-to-CVD pairs, $p = 0.668$, rank-biserial effect size and exact test statistic in Supplementary~\ref{app:effect_sizes}).
The same null held at the within-ROI level at hV4 (Crawford \& Howell $t$-test, both CVD subjects within the HC range, $p = 0.142$).
The 8 scanner stimuli are therefore individually distinguishable at the cortical level in both CVD participants, in line with prior work on CVD above-threshold identification \cite{bosten2019, boehm2014}.
The representational substrate on which a corrective filter could act is present.

% ──────────────────────────────────────────────────────────────────────────────
\subsection{Color interpolation is impaired in both CVD cases}
\label{sec:results:loco}
% ──────────────────────────────────────────────────────────────────────────────

We next localized where continuous hue interpolation fails, which defines the functional target of the corrective transform, using leave-one-color-out (LOCO) decoding of the forward encoding model.
Among the four retinotopic ROIs, only hV4 supported above-chance hue interpolation in HC (adjacent accuracy $0.47 \pm 0.05$, mean $\pm$ SEM, $p = 0.044$ under $8! = 40{,}320$ exact label permutations).
V1--V3 did not exceed chance despite preserving cross-run discrimination.
This confirms hV4 as the primary interpolation gate, replicating \citeA{brouwer2009}.

Both CVD participants fell at or below the adjacent-accuracy chance level at hV4 (Figure~\ref{fig:loco}B).
Sub-08 reached adjacent accuracy $0.25$ (Crawford \& Howell $t = -1.58$, $p = 0.082$, n.s., $d_{cc} = -1.71$).
Sub-09 reached $0.13$ ($t = -2.48$, $p = 0.024$, $d_{cc} = -2.68$).
The impairment was not uniform across hues.
Both participants showed near-zero adjacent accuracy specifically for S-cone intermediate colors at hV4, namely blue, purple, and magenta (Figure~\ref{fig:loco}C).
Per-hue Crawford \& Howell tests reached significance for blue (sub-08 $d = 2.13$, $p = 0.047$; sub-09 $d = 2.15$, $p = 0.046$) and for purple (sub-08 $d = 2.40$, $p = 0.033$).
These per-hue accuracies form each participant's hue vulnerability profile $\mathbf{v} \in [0,1]^8$.

% ── Figure 3 ─────────────────────────────────────────────────────────────────
\begin{figure*}[tb]
  \centering
  \includegraphics[width=\textwidth]{fig2_loro_loco}
  \caption{\textbf{Color discrimination is preserved but interpolation is selectively impaired at hV4 in CVD.} \textbf{(A)} LORO discrimination accuracy (forward encoding model, 8-class, SRM-aligned), V1--hV4.
  HC mean $\pm$ SEM ($n = 7$, gray); sub-08 (orange square), sub-09 (teal triangle).
  Chance = $1/8$.
  \textbf{(B)} LOCO adjacent accuracy ($\pm 1$ hue step; chance $= 3/8$).
  \emph{$*$}: sub-09 below HC at hV4 (Crawford \& Howell $t = -2.48$, $p = 0.024$; exploratory single-case finding); sub-08 $t = -1.58$, $p = 0.082$, n.s.
  \textbf{(C)} Per-hue adjacent accuracy at hV4, HC bars $\pm$ SEM.
  Per-hue Crawford \& Howell tests (one-tailed): blue --- sub-08 $d = 2.13$, $p = 0.047$; sub-09 $d = 2.15$, $p = 0.046$; purple --- sub-08 $d = 2.40$, $p = 0.033$.
  Joint pattern in (A) and (B--C) defines the filter precondition and target, respectively.}
  \label{fig:loco}
\end{figure*}

% ──────────────────────────────────────────────────────────────────────────────
\subsection{Color geometry is distorted at a distinct ROI in each CVD case}
\label{sec:results:geometry}
% ──────────────────────────────────────────────────────────────────────────────

To characterize the geometric basis of the LOCO impairment, we compared the pairwise representational structure of CVD and HC participants in a shared SRM-aligned space.
For each CVD participant we computed disparity and the per-pair RDM deviation from the HC mean ($\Delta\text{RDM} = \text{RDM}_\text{CVD} - \overline{\text{RDM}}_\text{HC}$).

The ROI specificity of elevated disparity diverges between the two CVD participants, with V2-dominant disparity in deutan and V1-dominant disparity in protan (Figure~\ref{fig:geometry}B).
Sub-08 showed significantly elevated disparity specifically at V2 (Crawford \& Howell $p = 0.040$), with no significant elevation at V1, V3, or hV4.
Sub-09 showed significantly elevated disparity specifically at V1 ($p = 0.007$), with no significant elevation at V2, V3, or hV4.
Full test statistics and effect sizes ($t$, $d_{cc}$) for these Crawford--Howell comparisons are reported in Supplementary~\ref{app:effect_sizes}.

The $\Delta$RDM heatmaps (Figure~\ref{fig:geometry}A) reveal subject-specific distortion structures.
Sub-08 shows elevated distances among S-cone intermediate pairs at V2, and sub-09 shows a broader reorganization of pairwise distances at V1.
The ROIs of strongest disparity elevation (V1 for sub-09, V2 for sub-08) differ from the ROI used as the distortion-model loss target (hV4, see the \nameref{sec:methods:twocomp} section). Mean pairwise disparity is an existence statement about absolute geometric perturbation in early visual cortex. hV4 is the cortical stage at which the corrective filter must act, so the $\Delta$RDM enters the loss as a geometric convergence check at that ROI. V1/V2 disparity at upstream ROIs is parallel evidence of distortion rather than a competing fit target.

% ── Figure 4 ─────────────────────────────────────────────────────────────────
\begin{figure*}[tb]
  \centering
  \includegraphics[width=\textwidth]{fig3_geometry}
  \caption{\textbf{Each CVD subject shows significantly elevated color representation disparity at a distinct ROI.} \textbf{(A)} $\Delta$RDM $= \text{RDM}_\text{CVD} - \overline{\text{RDM}}_\text{HC}$ in SRM-aligned space at each subject's primary ROI (sub-08 V2, sub-09 V1).
  Warm: CVD distances larger than HC mean; cool: smaller.
  \textbf{(B)} Mean pairwise correlation distance (disparity) per subject per ROI.
  HC band: group mean $\pm 1$ SD ($n = 7$); individual HC LOO as dots.
  Crawford \& Howell tests (sub-08 V2: $p = 0.040^*$; sub-09 V1: $p = 0.007^{**}$, exploratory single-case).
  (A) shows the directional pattern of pairwise deviations. (B) tests the absolute amplitude of pairwise distances against the HC LOO distribution.
  Idiosyncratic ROI specificity (V2 deutan; V1 protan) is inconsistent with a shared group-level gain mechanism.}
  \label{fig:geometry}
\end{figure*}

% ──────────────────────────────────────────────────────────────────────────────
\subsection{The retinal-plus-gain model is structurally insufficient}
\label{sec:results:rc_insufficient}
% ──────────────────────────────────────────────────────────────────────────────

The R+C model is structurally insufficient for both CVD participants.
Fitting this model to each CVD participant yielded ceiling saturation in
both cases.
For sub-08, $\langle 100\%\rangle$ of resample solutions saturated the
grid boundary ($g = \langle 3.0\rangle$).
For sub-09, saturation remained high ($\langle 41\%\rangle$,
$g = \langle 2.95\rangle$).
Boundary saturation at these rates indicates model misspecification,
not a well-constrained optimum \cite{wilson2019}.

A gain $g > 2$ implies that cortical processing reverses the direction of
the retinal shift past the undistorted hue
(Section~\ref{sec:methods:candidates}).
Both participants scored below the Ishihara criterion (sub-08: 5/14 plates
correct; sub-09: 7/14), confirming residual CVD under standard viewing
conditions.
A fitted $g > 2$ is therefore internally inconsistent with confirmed CVD
status: the model claims overcompensation where the participants
demonstrably show impairment.
The saturated boundary is a model failure, not a valid parameter estimate.

The structural cause is a degree-of-freedom deficit.
The function $\delta\theta = (2 - g)\cdot\delta\theta_{\rm Mach}$
displaces hues only along the fixed confusion axis.
No value of $g$ can account for distortions in directions orthogonal to
that axis.
R+C is not considered further as a filter candidate.

% \textcolor{red}{[TODO: R2 --- A single cortical model fits both CVD cases]}
% \textcolor{red}{[TODO: R2.5 --- Neural data identifies what behavior cannot]}
% \textcolor{red}{[TODO: R3 --- Identifiability: mechanism class robust, per-axis value not]}
% \textcolor{red}{[TODO: Figure 4 --- loss landscape caption (6,-42)/(2,+24) 기준으로 교체]}

% ──────────────────────────────────────────────────────────────────────────────
\subsection{\textcolor{red}{[TODO] Per-subject stimulus-space filter}}
\label{sec:results:filter}
% ──────────────────────────────────────────────────────────────────────────────

{\color{red}
The 2-component transform is bijective, so a stimulus-space pre-image $\tilde{\theta}_k$ such that $T(\tilde{\theta}_k) = \theta_k$ exists for every displayed hue $\theta_k$ (see the \nameref{sec:methods:filter} section).
The per-subject correction filter is $\delta\theta^{\rm filt}_k = \tilde{\theta}_k - \theta_k$.
For both participants, 8/8 pre-images were exact (residual $< 0.001^\circ$, Figure~\ref{fig:filter}).

For sub-08 (deutan, $\hat\beta_s = 38^\circ$, $\hat\beta_c = -14^\circ$, Stockman confusion axis $\theta_{\rm conf} = 150^\circ$), the filter applies a wide-amplitude L/M-rotation.
A near-uniform negative shift is applied across red, orange, yellow, green, and cyan ($\delta\theta^{\rm filt}$ ranging $-10^\circ$ to $-32^\circ$, peak $-32^\circ$ at cyan, $\theta = 180^\circ$).
A large opposite-sign rotation is applied at purple ($+29^\circ$, $\theta = 270^\circ$) and magenta ($+19^\circ$, $\theta = 315^\circ$).
The mean correction magnitude is $|\delta\theta^{\rm filt}| = 22^\circ$.

For sub-09 (protan, $\hat\beta_s = 6^\circ$, $\hat\beta_c = -22^\circ$, Stockman axis $\theta_{\rm conf} = 16^\circ$), the correction is a far smaller monotonic rotation along the protan confusion axis.
Peak displacements are at blue ($+21^\circ$, $\theta = 225^\circ$) and magenta ($-21^\circ$, $\theta = 315^\circ$), with near-zero shifts at green and cyan.
The mean correction magnitude is $|\delta\theta^{\rm filt}| = 10^\circ$.

The two correction signatures reflect distinct CVD etiologies rather than mirrored protan/deutan symmetry.
Under the R+C decomposition (Appendix~\ref{app:altmodels}), sub-08 is cortical-dominant and sub-09 is retinal-dominant.
A single retinal-family filter cannot serve both individuals.

The model-class expressiveness has a closed-form upper bound.
The magnitude $|\delta\theta|$ at $\theta = 45^\circ$ is bounded by $|0.71\,\hat\beta_s - 0.26\,\hat\beta_c|$.
Corrective shifts exceeding $\sim\!51^\circ$ at the deutan mid-axis (e.g.\ an orange-to-green excursion of $\gtrsim 85^\circ$) lie outside the 2-component class for any biologically plausible $\hat\beta$ and would require an additional mechanism such as saturation or lightness.
}

% ── Figure 6 ─────────────────────────────────────────────────────────────────
\begin{figure*}[!htbp]
  \centering
  \includegraphics[width=\textwidth,height=0.85\textheight,keepaspectratio]{fig5_preimage}
  \caption[Per-subject 4-column filter rendering]{\textbf{Per-subject stimulus-space filter: 4-column rendering, side-by-side.}
  Left block: sub-08 (deutan), filter at $(\hat\beta_s, \hat\beta_c) = (38^\circ, -14^\circ)$, $\|\hat\beta\| = 40.5^\circ$.
  Right block: sub-09 (protan), filter at $(\hat\beta_s, \hat\beta_c) = (6^\circ, -22^\circ)$, $\|\hat\beta\| = 22.8^\circ$.
  Eight rows correspond to the eight displayed DKL hues.
  Columns (left to right) within each block:
  \emph{Original} (HC percept of the displayed stimulus);
  \emph{CVD percept} (simulated CVD percept of the original under the per-subject 2-component model);
  \emph{Filtered} (stimulus after applying $\delta\theta^{\rm filt}$, the exact pre-image of $-\delta\theta(\hat\beta_s, \hat\beta_c)$ computed by Brent's-method refinement of the forward map; 8/8 exact, residual $< 0.001^\circ$ for both subjects);
  \emph{CVD(Filt.)} (simulated CVD percept of the filtered stimulus).
  Rendering coordinate space: STIM\_LAB CIELab (project convention).
  Quantitative behavioral validation of filter efficacy is deferred to the Phase-3 2AFC arm and will be added once data are collected.}
  \label{fig:filter}
\end{figure*}

% ──────────────────────────────────────────────────────────────────────────────
\subsection{\textcolor{red}{[TODO] Behavioral and neural filter evaluation}}
\label{sec:results:filter_eval}
% ──────────────────────────────────────────────────────────────────────────────

\textcolor{red}{Filter efficacy will be evaluated at two levels. (1)~Behavioral: same-day 2-alternative forced-choice (2AFC) discrimination of pre-specified hue pairs under the per-subject filter versus a no-filter baseline. (2)~Neural: second-session fMRI under the 2-component-derived filter versus a generic Windows accessibility filter (counterbalanced), with LOCO and LORO decoded on second-session data (see the \nameref{sec:methods:behavioral} section). Results to be added once data analysis is complete.}
